% Figure: a Riemann sum with left endpoints, right endpoints, and the
% true area under the curve f(x) = 1 + 0.6 sin(2x) on [0, pi].
\begin{figure}[htbp]
\centering
\begin{tikzpicture}
\begin{axis}[
  width=12cm, height=6cm,
  xlabel={$x$},
  ylabel={$f(x)$},
  xmin=0, xmax=3.5,
  ymin=0, ymax=2.2,
  domain=0:3.1416,
  samples=200,
  axis lines=left,
  grid=both,
  major grid style={engblack!15},
  font=\small,
  legend pos=south west,
  legend cell align=left,
]

% True area under the curve, shaded
\addplot[draw=none, fill=engblue!12, domain=0:3.1416]
  {1 + 0.6*sin(deg(2*x))} \closedcycle;

% Left-endpoint rectangles, n = 6 on [0, pi]
\foreach \i in {0,1,2,3,4,5}{
  \pgfmathsetmacro{\xl}{\i*3.1416/6}
  \pgfmathsetmacro{\xr}{(\i+1)*3.1416/6}
  \pgfmathsetmacro{\yl}{1 + 0.6*sin(deg(2*\xl))}
  \addplot[draw=engred, fill=engred!10, fill opacity=0.4]
    coordinates {(\xl,0) (\xr,0) (\xr,\yl) (\xl,\yl)} \closedcycle;
}

% The curve itself
\addplot[engblack, very thick] {1 + 0.6*sin(deg(2*x))};

\legend{area $\int_{0}^{\pi} f$,
        left-endpoint rectangles ($n=6$),
        $f(x) = 1 + 0.6\sin(2x)$}

\end{axis}
\end{tikzpicture}
\caption[Riemann sum as area accounting]{The Riemann sum as area
accounting. The shaded region under $f(x) = 1 + 0.6\sin(2x)$ from
$0$ to $\pi$ is the definite integral. The red rectangles, with
height taken at the left endpoint of each subinterval and width
$\pi/6$, form one Riemann sum that overshoots on rising portions of
$f$ and undershoots on falling portions. As the mesh of the
partition tends to zero, every choice of sample points and every
partition gives the same limit, which is the integral.}
\label{fig:vol02:ch06:riemann-sum}
\end{figure}
